\subsubsection{}

Argon-atom 1 har posisjon $(0, 0)$ og argon-atom 2 har posisjon $(0, y)$.
Avstanden mellom atomene er da $r = y$.

Som vi har sett tidligere er den potensielle energien integralet av kraften som virker på partikkelen. Motsatt er altså kraften den deriverte av energien.

Kraften mellom atomene er gitt ved den negative gradienten av potensialet.
I polarkoordinater med $\hat{e}_r$ pekende fra atom 1 mot atom 2:

\begin{align}
F(r) = -\frac{dE_p}{dr} = C\left(12\frac{l_0^{12}}{r^{13}} - 6\frac{l_0^6}{r^7}\right)
\end{align}

Med $C = 1.735 \times 10^{-21}$\,J og $l_0 = 2.36 \times 10^{-10}$\,m og $r = y$:

\begin{align}
F(y) = 1.735 \times 10^{-21} \left(12\frac{(2.36 \times 10^{-10})^{12}}{y^{13}} - 6\frac{(2.36 \times 10^{-10})^{6}}{y^{7}}\right)
\end{align}

For kraften $\vec{F}_{1\to2}$ peker $\hat{e}_r$ fra atom 1 mot atom 2,
altså i retning $(0, y)$, dvs.\ i positiv $y$-retning $\hat{j}$:

\begin{align}
\vec{F}_{1\to2} = F(y)\,\hat{j}
\end{align}

For kraften $\vec{F}_{2\to1}$ peker $\hat{e}_r$ fra atom 2 mot atom 1,
altså i negativ $y$-retning $-\hat{j}$:

\begin{align}
\vec{F}_{2\to1} = -F(y)\,\hat{j}
\end{align}

Vi ser at $\vec{F}_{2\to1} = -\vec{F}_{1\to2}$, altså er kreftene like store og motsatt rettet.
Newtons 3. lov er dermed overholdt.

\newpage